\begin{theorem}[正则黄金二次层与 Lee--Yang 奇性三次层的最小正规公共载体]\label{thm:conclusion-leyang-global-local-minimal-normal-closure}
沿用
\[
K_{\mathrm{reg}}:=J_{\mathrm{pert}}=\QQ(\sqrt5),
\qquad
K_{\mathrm{sing}}:=K_{\mathrm{br}}=\QQ(u),
\]
并记 \(H_{\mathrm{sing}}\) 为 \(K_{\mathrm{sing}}\) 的 Galois 闭包，设
\[
\Omega:=K_{\mathrm{reg}}H_{\mathrm{sing}}.
\]
则
\[
[K_{\mathrm{reg}}K_{\mathrm{sing}}:\QQ]=6,
\qquad
[\Omega:\QQ]=12,
\qquad
\Gal(\Omega/\QQ)\cong C_2\times S_3.
\]
并且 \(\Omega\) 正是 \(K_{\mathrm{reg}}K_{\mathrm{sing}}\) 的最小正规闭包。
\end{theorem}
\begin{proof}
由定理 \ref{thm:conclusion-local-cumulant-golden-quadratic-leyang-cubic-disjointness}，
\[
K_{\mathrm{reg}}\cap K_{\mathrm{sing}}=\QQ,
\qquad
[K_{\mathrm{reg}}K_{\mathrm{sing}}:\QQ]=2\cdot 3=6.
\]
再由定理 \ref{thm:conclusion-leyang-dominant-edge-quadratic-resolvent-minus111} 与命题 \ref{prop:conclusion-leyang-abelian-shadow-nonabelian-core}，\(H_{\mathrm{sing}}/\QQ\) 的 Galois 群为 \(S_3\)，且其唯一二次子域为
\[
K_{\mathrm{LY}}=\QQ(\sqrt{-111}).
\]
由于
\[
K_{\mathrm{reg}}=\QQ(\sqrt5)\neq \QQ(\sqrt{-111})=K_{\mathrm{LY}},
\]
故
\[
K_{\mathrm{reg}}\cap H_{\mathrm{sing}}=\QQ.
\]
而 \(K_{\mathrm{reg}}/\QQ\) 与 \(H_{\mathrm{sing}}/\QQ\) 都是 Galois 扩张，因此
\[
[\Omega:\QQ]=[K_{\mathrm{reg}}:\QQ]\,[H_{\mathrm{sing}}:\QQ]=2\cdot 6=12,
\]
且
\[
\Gal(\Omega/\QQ)\cong \Gal(K_{\mathrm{reg}}/\QQ)\times \Gal(H_{\mathrm{sing}}/\QQ)\cong C_2\times S_3.
\]
由于 \(H_{\mathrm{sing}}\) 本身就是 \(K_{\mathrm{sing}}\) 的最小正规闭包，而 \(K_{\mathrm{reg}}\) 已经正规，故 \(\Omega\) 正是 \(K_{\mathrm{reg}}K_{\mathrm{sing}}\) 的最小正规闭包。证毕。
\end{proof}

\begin{corollary}[最大全阿贝尔阴影与被迫出现的第三二次子域]\label{cor:conclusion-leyang-global-local-maximal-abelian-shadow}
\(\Omega/\QQ\) 的极大阿贝尔子扩张恰为
\[
\Omega^{\mathrm{ab}}=\QQ(\sqrt5,\sqrt{-111}),
\]
因此它是一个双二次域，其三个且仅有三个二次子域为
\[
\QQ(\sqrt5),\qquad
\QQ(\sqrt{-111}),\qquad
\QQ(\sqrt{-555}).
\]
\end{corollary}
\begin{proof}
由定理 \ref{thm:conclusion-leyang-global-local-minimal-normal-closure}，
\[
\Gal(\Omega/\QQ)\cong C_2\times S_3.
\]
其交换化为
\[
(C_2\times S_3)^{\mathrm{ab}}\cong C_2\times C_2,
\]
而交换子群正是 \(A_3\subset S_3\)。故极大阿贝尔子扩张为固定域 \(\Omega^{A_3}\)，其次数为 \(4\)。

另一方面，\(\Omega\) 同时包含 \(K_{\mathrm{reg}}=\QQ(\sqrt5)\) 与 \(H_{\mathrm{sing}}\) 的唯一二次子域 \(\QQ(\sqrt{-111})\)，因而
\[
\QQ(\sqrt5,\sqrt{-111})\subseteq \Omega^{A_3}.
\]
左端已经是次数 \(4\) 的双二次域，故必有
\[
\Omega^{\mathrm{ab}}=\Omega^{A_3}=\QQ(\sqrt5,\sqrt{-111}).
\]
双二次域的第三个二次子域由平方类乘积给出，即
\[
\QQ(\sqrt5)\cdot \QQ(\sqrt{-111})=\QQ(\sqrt{-555}).
\]
证毕。
\end{proof}

\begin{theorem}[阿贝尔化必然抹去 Lee--Yang 奇性分支的三张 sheet]\label{thm:conclusion-leyang-global-local-abelianization-kills-sheets}\leanverified{paper\_conclusion\_leyang\_global\_local\_abelianization\_kills\_sheets}
对任意阿贝尔群 \(A\) 与任意群同态
\[
\rho:\Gal(\Omega/\QQ)\longrightarrow A,
\]
\(\rho\) 都必湮灭 \(A_3\subset S_3\)。从而，一切阿贝尔可观测量都无法区分 \(f(T)=324T^3-540T^2+333T-74\) 的三根所对应的三张奇性 sheet。等价地，所有阿贝尔观测都只能下降到结论 \ref{cor:conclusion-leyang-global-local-maximal-abelian-shadow} 的双二次阴影，而绝不可能恢复 \(u\) 本身。
\end{theorem}
\begin{proof}
任意到阿贝尔群的群同态都必经交换化，故必湮灭交换子群
\[
[\Gal(\Omega/\QQ),\Gal(\Omega/\QQ)]=A_3.
\]
但在 \(H_{\mathrm{sing}}\) 中，\(A_3\) 正是循环置换三根的三循环子群。因此经由任意阿贝尔观测，三张奇性 sheet 必落入同一不可区分轨道。

若 \(u\) 可由某个阿贝尔观测恢复，则必有
\[
u\in \Omega^{\mathrm{ab}}.
\]
然而推论 \ref{cor:conclusion-leyang-branch-coordinate-no-finite-abelian-ledger} 已证明：\(u\) 不属于任何有限阿贝尔扩张。由于 \(\Omega^{\mathrm{ab}}/\QQ\) 是有限阿贝尔扩张，矛盾。故 \(u\notin \Omega^{\mathrm{ab}}\)，阿贝尔观测不可能恢复 \(u\)。证毕。
\end{proof}

\begin{theorem}[Lee--Yang 奇性三次数域在 \texorpdfstring{$5$}{5} 处的无分歧循环化]\label{thm:conclusion-leyang-global-local-q5-unramified-cubic}
\leanverified{paper\_conclusion\_leyang\_global\_local\_q5\_unramified\_cubic}
记
\[
K_{\mathrm{sing},5}:=K_{\mathrm{sing}}\otimes_{\QQ}\QQ_5.
\]
则 \(K_{\mathrm{sing},5}\) 是 \(\QQ_5\) 的唯一三次无分歧扩张。特别地，
\[
[K_{\mathrm{sing},5}:\QQ_5]=3,
\qquad
e(K_{\mathrm{sing},5}/\QQ_5)=1,
\qquad
f(K_{\mathrm{sing},5}/\QQ_5)=3,
\]
并且
\[
\Gal(K_{\mathrm{sing},5}/\QQ_5)\cong C_3.
\]
\end{theorem}
\begin{proof}
把
\[
f(T)=324T^3-540T^2+333T-74
\]
模 \(5\) 约化，得到
\[
\bar f(T)\equiv 4T^3+3T+1.
\]
直接代入 \(T=0,1,2,3,4\) 可见 \(\bar f\) 在 \(\FF_5\) 中无根。由于 \(\deg \bar f=3\)，故 \(\bar f\) 在 \(\FF_5[T]\) 中不可约。

另一方面，由定理 \ref{thm:conclusion-leyang-dominant-edge-quadratic-resolvent-minus111} 的证明，\(f\) 的判别式为
\[
\Disc(f)=-2^6\cdot 3^9\cdot 37,
\]
故 \(5\nmid \Disc(f)\)。因此 \(5\) 在 \(K_{\mathrm{sing}}\) 中不分歧；又因 \(\bar f\) 不可约，\(5\) 对应唯一素理想，且剩余次数为 \(3\)。于是
\[
e=1,\qquad f=3.
\]
对局部域而言，无分歧三次扩张必为 Galois 且其 Galois 群循环，故
\[
\Gal(K_{\mathrm{sing},5}/\QQ_5)\cong C_3.
\]
证毕。
\end{proof}

\begin{corollary}[正则--奇性联合 \texorpdfstring{$5$}{5}-进载体是 \texorpdfstring{$(e,f)=(2,3)$}{(e,f)=(2,3)} 的循环六次扩张]\label{cor:conclusion-leyang-global-local-q5-composite-c6}
记
\[
L_5:=\QQ_5(\sqrt5,u).
\]
则
\[
[L_5:\QQ_5]=6,
\qquad
e(L_5/\QQ_5)=2,
\qquad
f(L_5/\QQ_5)=3,
\]
并且
\[
\Gal(L_5/\QQ_5)\cong C_6.
\]
\end{corollary}
\begin{proof}
由推论 \ref{cor:conclusion-fibadic-golden-extension-even-ramification-necessity}，
\[
\QQ_5(\sqrt5)/\QQ_5
\]
是二次全分歧扩张。由定理 \ref{thm:conclusion-leyang-global-local-q5-unramified-cubic}，
\[
K_{\mathrm{sing},5}/\QQ_5
\]
是三次无分歧扩张。无分歧扩张与全分歧扩张交域必为基域，故
\[
\QQ_5(\sqrt5)\cap K_{\mathrm{sing},5}=\QQ_5.
\]
于是
\[
[L_5:\QQ_5]=2\cdot 3=6,
\qquad
e(L_5/\QQ_5)=2,
\qquad
f(L_5/\QQ_5)=3.
\]
两侧子扩张都为 Galois，且 Galois 群分别为 \(C_2\) 与 \(C_3\)；因此合成扩张也为 Galois，且
\[
\Gal(L_5/\QQ_5)\cong C_2\times C_3\cong C_6.
\]
证毕。
\end{proof}

\begin{theorem}[全局非阿贝尔性在 \texorpdfstring{$5$}{5}-进纤维上完全隐没]\label{thm:conclusion-leyang-global-local-nonabelianity-disappears-at-q5}
设 \(\mathfrak P\) 为 \(\Omega\) 上位于 \(5\) 之上的素理想。则其分解群满足
\[
D_{\mathfrak P}\cong C_6.
\]
特别地，
\[
\Gal(\Omega/\QQ)\cong C_2\times S_3
\]
是非阿贝尔群，而所有 \(\QQ_5\)-局部完成所见到的分解群却都是阿贝尔的。故 Lee--Yang 奇性分支的非阿贝尔性并不诞生于 \(5\)-进局部纤维，而只诞生于全局拼接。
\end{theorem}
\begin{proof}
由定理 \ref{thm:conclusion-leyang-global-local-q5-unramified-cubic}，\(K_{\mathrm{sing},5}/\QQ_5\) 已经是 Galois 的循环三次扩张。因此 \(H_{\mathrm{sing}}\) 在 \(5\) 之上的局部完成不会比 \(K_{\mathrm{sing},5}\) 更大；换言之，\(H_{\mathrm{sing}}\) 的任一 \(5\)-进完成都同构于 \(K_{\mathrm{sing},5}\)。

再与 \(\QQ_5(\sqrt5)\) 合成，便得到结论 \ref{cor:conclusion-leyang-global-local-q5-composite-c6} 中的 \(L_5\)。因此
\[
\Omega_{\mathfrak P}\cong L_5.
\]
分解群与局部 Galois 群同构，故
\[
D_{\mathfrak P}\cong \Gal(\Omega_{\mathfrak P}/\QQ_5)\cong \Gal(L_5/\QQ_5)\cong C_6.
\]
这说明 \(S_3\)-型非交换性在 \(5\)-进纤维上完全退化为循环六次局部单调性。证毕。
\end{proof}

\begin{theorem}[Lee--Yang 振幅常数的 \texorpdfstring{$5$}{5}-进提升只增加剩余次数而不增加分歧指数]\label{thm:conclusion-leyang-global-local-amplitude-b-q5-lift}
设 \(B\) 为 Lee--Yang 振幅常数，并记
\[
L_B:=\QQ(B).
\]
则 \(B\) 的最小多项式可取为
\[
h(X)=162X^6+1593X^4+1998X^2+1184.
\]
并且
\[
L_{B,5}:=L_B\otimes_{\QQ}\QQ_5
\]
是 \(\QQ_5\) 的唯一六次无分歧扩张；因此
\[
[L_{B,5}:\QQ_5]=6,
\qquad
e(L_{B,5}/\QQ_5)=1,
\qquad
f(L_{B,5}/\QQ_5)=6.
\]
进一步，
\[
[\QQ(\sqrt5,B):\QQ]=12,
\]
且其 \(\QQ_5\)-完成满足
\[
[\QQ_5(\sqrt5,B):\QQ_5]=12,
\qquad
e=2,
\qquad
f=6,
\]
并有
\[
\Gal(\QQ_5(\sqrt5,B)/\QQ_5)\cong C_2\times C_6.
\]
\end{theorem}
\begin{proof}
由定理 \ref{thm:xi-terminal-zm-stokes-amplitude-degree6-minpoly}，
\[
[\QQ(B):\QQ]=6,
\]
且 \(B\) 的极小多项式正是
\[
h(X)=162X^6+1593X^4+1998X^2+1184.
\]
把 \(h\) 模 \(5\) 约化，得
\[
\bar h(X)\equiv 2X^6-2X^4-2X^2-1.
\]
直接在 \(\FF_5[X]\) 中计算可得
\[
\gcd(\bar h,X^5-X)=\gcd(\bar h,X^{25}-X)=\gcd(\bar h,X^{125}-X)=1.
\]
因此 \(\bar h\) 没有 \(1\)、\(2\)、\(3\) 次因子；而 \(\deg \bar h=6\)，遂 \(\bar h\) 在 \(\FF_5[X]\) 中不可约。

再由定理 \ref{thm:xi-leyang-amplitude-B-sextic-biquadratic-shadow}，
\[
\Disc(h)=-2^{16}\cdot 3^{22}\cdot 11^4\cdot 37^3\cdot 109^4,
\]
故 \(5\nmid \Disc(h)\)。于是 \(L_{B,5}/\QQ_5\) 是六次无分歧扩张，从而必为唯一六次无分歧扩张。

另一方面，\(\QQ_5(\sqrt5)/\QQ_5\) 是二次全分歧扩张，而 \(L_{B,5}/\QQ_5\) 无分歧，因此 \(L_{B,5}\) 不可能包含 \(\QQ_5(\sqrt5)\)。于是
\[
\QQ(B)\cap \QQ(\sqrt5)=\QQ,
\]
从而
\[
[\QQ(\sqrt5,B):\QQ]=2\cdot 6=12.
\]
局部上，无分歧六次扩张与全分歧二次扩张交域平凡，故
\[
[\QQ_5(\sqrt5,B):\QQ_5]=12,
\qquad
e=2,
\qquad
f=6.
\]
两侧局部子扩张都为 Galois，因此其合成扩张也为 Galois，且 Galois 群为
\[
C_2\times C_6.
\]
证毕。
\end{proof}

\begin{theorem}[两切片最小完备阈值与离开有限阿贝尔算术账本的阈值重合]\label{thm:conclusion-leyang-global-local-two-slice-threshold}
在单原子差分通道
\[
c^{(j)}=m\tau_j^E
\qquad (j=1,2)
\]
中，单切片永远不足以唯一恢复参数对 \((m,E)\)，而任意两条不同切片都足以通过
\[
E=\frac{\log(c^{(2)}/c^{(1)})}{\log(\tau_2/\tau_1)},
\qquad
m=c^{(1)}\tau_1^{-E}
\]
精确恢复 \((m,E)\)。当指数参数特化为 Lee--Yang 奇性坐标
\[
E=u=\kappa_\star^2
\]
时，最小解析完备阈值与离开有限阿贝尔算术账本的阈值严格重合：达到“两切片”后即可精确恢复 \(u\)，而所恢复出的 \(u\) 立刻不属于任何有限阿贝尔扩张。
\end{theorem}
\begin{proof}
前文已经证明：单切片不足以唯一恢复 \((m,E)\)，而任意两条不同切片则给出上述显式重构公式。把 \(E\) 特化为 \(u\) 后，便得到
\[
u=\frac{\log(c^{(2)}/c^{(1)})}{\log(\tau_2/\tau_1)}.
\]
因此，两切片已经足以精确恢复奇性坐标本身。

另一方面，由推论 \ref{cor:conclusion-leyang-branch-coordinate-no-finite-abelian-ledger}，\(u\) 不属于任何有限阿贝尔扩张。故一旦跨过“恰好两切片”的最小解析完备阈值，所得参数立即越出全部有限阿贝尔 residue 账本。证毕。
\end{proof}
